\documentclass{article}

% Math packages (fixes \mathbb, \dagger, \longmapsto, etc.)
\usepackage{amsmath, amssymb, amsfonts}

\title{Matrix Lie Algebras II - Gauge Theory 2025-2026}
\author{Soham Kandalgaonkar}
\date{February 2026}
\usepackage{amsmath,amssymb,amsthm}

\theoremstyle{definition}
\newtheorem{definition}{Definition}[section]
\begin{document}

\maketitle

\section{What is $GL(n,\mathbb{C})$?}

\begin{definition}
The \textbf{general linear group of degree $n$ over $\mathbb{C}$}, denoted $GL(n,\mathbb{C})$, is defined as
\[
GL(n,\mathbb{C}) = \{ A \in M_n(\mathbb{C}) \mid \det(A) \neq 0 \}.
\]
\end{definition}

Equivalently, $GL(n,\mathbb{C})$ is the set of all invertible $n \times n$ matrices with complex entries. Since a matrix is invertible if and only if its determinant is nonzero, the two characterizations agree.


\subsection*{Geometric Structure}

The space $M_n(\mathbb{C})$ of all complex $n \times n$ matrices can be identified with
\[
\mathbb{C}^{n^2} \cong \mathbb{R}^{2n^2}.
\]

The determinant
\[
\det : M_n(\mathbb{C}) \to \mathbb{C}
\]
is a polynomial map in the matrix entries and, therefore, continuous.

Since
\[
GL(n,\mathbb{C}) = \det^{-1}(\mathbb{C} \setminus \{0\}),
\]
and $\mathbb{C} \setminus \{0\}$ is an open subset of $\mathbb{C}$, it follows that
\[
GL(n,\mathbb{C})
\]
is an open subset of $M_n(\mathbb{C})$.


Therefore,

\[
GL(n,\mathbb{C})
\text{ is a smooth real manifold of dimension } 2n^2.
\]

\subsection*{Lie Algebra of $GL(n,\mathbb{C})$}

The Lie algebra of $GL(n,\mathbb{C})$ is

\[
\mathfrak{gl}(n,\mathbb{C}) = M_n(\mathbb{C})
\]

with Lie bracket defined by the commutator:

\[
[X,Y] = XY - YX.
\]

The tangent space at the identity satisfies the following condition.

\[
T_I GL(n,\mathbb{C}) = M_n(\mathbb{C}),
\]

\textit{Proof.}
Let $M_n(\mathbb{C})$ be a real vector space of dimension $2n^2$ and note that
$GL(n,\mathbb{C}) \subset M_n(\mathbb{C})$. We first must show that
$GL(n,\mathbb{C})$ is an open subset of $M_n(\mathbb{C})$.
The determinant map
\[
\det : M_n(\mathbb{C}) \to \mathbb{C}
\]
is continuous, as it is a polynomial of its matrix entries. Previously stated,
\[
GL(n,\mathbb{C})=\det^{-1}(\mathbb{C}\setminus\{0\}).
\]
Since $\mathbb{C} \setminus \{0\}$ is open in $\mathbb{C}$, it follows that
$GL(n,\mathbb{C})$ is open in $M_n(\mathbb{C})$. Now let there be
finite-dimensional real vector space $V$ such that
\[
U \subset V.
\]
Then for any $p \in U$, the tangent space $T_pU$ is in $V$.

Now, let
\[
U = GL(n,\mathbb{C})
\]
\[
V = M_n(\mathbb{C})
\]
\[
p = I
\]

we obtain
\[
T_I GL(n,\mathbb{C}) \cong M_n(\mathbb{C}).
\]

With the definition of a tangent space of an open subset with a vector space itself, we get:
\[
T_I GL(n,\mathbb{C}) = M_n(\mathbb{C}).
\]

This proves our claim. \(\square\)

For every matrix $X \in M_n(\mathbb{C})$ generates a one-parameter subgroup

\[
A(t) = e^{tX}.
\]

Thus, $\mathfrak{gl}(n,\mathbb{C})$ consists of all complex $n \times n$ matrices equipped with the commutator bracket.

\[
\text{Lie Product Formula: }\text{For any } X,Y \in \mathfrak{gl}(n,\mathbb{C}),
\text{ the exponential of their sum can be expressed as:}
\]

\[
e^{X+Y} \;=\; \lim_{m\to\infty}\left(e^{X/m}e^{Y/m}\right)^m.
\]

\subsection{The Special Linear Group $SL(2,\mathbb{R})$ is defined as:}


\[
SL(2,\mathbb{R})=\left\{\,A\in M_2(\mathbb{R}) \;\middle|\; \det(A)=1\,\right\}
\]

\[
\text{Its associated Lie algebra } \mathfrak{sl}(2,\mathbb{R}) \text{ consists of all } 2\times 2
\text{ real matrices with zero trace:}
\]

\[
\mathfrak{sl}(2,\mathbb{R})=\left\{\,X\in M_2(\mathbb{R}) \;\middle|\; \operatorname{Tr}(X)=0\,\right\}
\]

\section{Properties of the Hermitian Conjugate}

Let $U$ be a complex $n \times n$ matrix.

\begin{definition}
The \textbf{Hermitian conjugate} (or \textbf{conjugate transpose}) of $U$, denoted $U^\dagger$, is defined by
\[
U^\dagger = \overline{U}^{\,T}.
\]
\end{definition}

That is, the Hermitian conjugate is obtained by:

\begin{itemize}
    \item Taking the transpose (interchanging rows and columns),
    \item Taking the complex conjugate of each entry.
\end{itemize}
\subsection{Standard Inner Product}
The standard inner product on $\mathbb{C}^n$ is:
\[
\langle a,b\rangle=\sum_{i=1}^n \overline{a_i}\,b_i.
\]
\[
\text{The standard inner product on } \mathbb{C}^n \text{ satisfies the positivity condition:}
\]

\[
\langle v, v\rangle \ge 0 \ \text{for all } v\in\mathbb{C}^n,
\qquad
\langle v, v\rangle = 0 \iff v=0.
\]

\[
\text{This allows us to define the norm of a vector as }
\|v\|=\sqrt{\langle v,v\rangle}.
\]

\subsection*{Entry-by-Entry Formula}

If $U = [u_{ij}]$, then

\[
(U^\dagger)_{ij} = \overline{u_{ji}}.
\]

\subsection*{Example}

Let

\[
U =
\begin{pmatrix}
2 + i & 2 \\
3i & 5 - i
\end{pmatrix}.
\]

Then

\[
U^\dagger =
\begin{pmatrix}
2 - i & -3i \\
2 & 5 + i
\end{pmatrix}.
\]

\subsection*{Algebraic Properties}

For complex matrices $A,B$ and scalar $\lambda \in \mathbb{C}$:

\begin{align*}
(A^\dagger)^\dagger &= A, \\
(A+B)^\dagger &= A^\dagger + B^\dagger, \\
(\lambda A)^\dagger &= \overline{\lambda}\, A^\dagger, \\
(AB)^\dagger &= B^\dagger A^\dagger.
\end{align*}
\subsection{Hermitian and Unitary Matrices}
A matrix \(H\) is Hermitian (or self-adjoint) if \(H = H^\dagger\).
An equivalent characterization using the inner product is:

\[
\langle v, Hw\rangle = \langle Hv, w\rangle
\quad \text{for all } v,w \in \mathbb{C}^n.
\]

\[
\text{A consequence for Hermitian matrices is that }
\langle v, Hv\rangle \in \mathbb{R}
\quad \text{for all } v.
\]
\subsection*{Unitary Matrices}

A matrix $U$ is \textbf{unitary} if

\[
U^\dagger U = I = U U^\dagger.
\]

Equivalently,

\[
U^{-1} = U^\dagger.
\]

\textit{Proof.}
A matrix $U$ is unitary precisely when it is adjoint to its inverse. By definition, $U$ is unitary if
\[
U^\dagger = U^{-1}.
\]
This is equivalent to saying that $U^\dagger$ is both a left inverse and a right inverse of $U$, that is,
\[
U^\dagger U = I
\qquad \text{and} \qquad
UU^\dagger = I.
\]
Hence, $U$ is unitary if
\[
U^\dagger U = I = UU^\dagger.
\]
\hfill $\square$

The unitary group is

\[
U(n) = \{ U \in GL(n,\mathbb{C}) \mid U^\dagger U = I \}.
\]

Additionally, 

\[
\text{A matrix } U \text{ is unitary if and only if it preserves the inner product:}
\]

\[
\langle v,w\rangle = \langle Uv, Uw\rangle
\quad \text{for all } v,w \in \mathbb{C}^n.
\]

\subsection*{Relation to Lie Algebra of $U(n)$}

Let $A(t)$ be a smooth path in $U(n)$ with $A(0) = I$.  
Since

\[
A(t) A^\dagger(t) = I,
\]

differentiating at $t=0$ gives

\[
A'(0) + A'(0)^\dagger = 0.
\]

Thus, the tangent space at the identity consists of matrices $X$ satisfying

\[
X^\dagger = -X.
\]

These matrices are \textbf{anti-Hermitian}.

Therefore, the Lie algebra of $U(n)$ is

\[
\mathfrak{U}(n)
=
\{ X \in M_n(\mathbb{C}) \mid X^\dagger = -X \}.
\]


\end{document}
